\documentclass[final]{beamer}
% you can chose your theme here:
\usetheme{I6dv}
% further beamerposter themes are available at
% http://www-i6.informatik.rwth-aachen.de/~dreuw/latexbeamerposter.php
  \usepackage{type1cm}
  \usepackage{calc}
  \usepackage{times}
  \usepackage{amsmath,amsthm, amssymb, latexsym}
\usepackage{tikz}
\usepackage{tikz-bpmn}
\usetikzlibrary{positioning}
\usepackage[utf8]{inputenc}  %or: \usepackage[latin1]{inputenc}
  \boldmath
  \usepackage[english]{babel}
  \usepackage[orientation=portrait,size=a0,scale=1.5]{beamerposter}
\setbeamersize{text margin left=18mm}
\setbeamersize{text margin right=18mm}
  \title{gLite sur Grid'5000~:\\[0.1em] vers une plate-forme d'expérimentation\\[0.1em] à taille réelle pour les grilles de production}
\author{Sébastien Badia et Lucas Nussbaum}
\institute{LORIA / Nancy-Université \hskip 2em INRIA Nancy -- Grand Est}

%  \newcommand{\footlinetext}{Lehrstuhl f\"ur Informatik 6 - Computer Science Department - RWTH Aachen University - Aachen, Germany \par Mail: \texttt{<surname>@cs.rwth-aachen.de} \hfill WWW: \texttt{http://www-i6.informatik.rwth-aachen.de}\vskip1ex}
%  \date{Jul. 31th, 2007}

\setbeamerfont{itemize/enumerate subbody}{size=\normalsize}
\setbeamerfont{itemize/enumerate subsubbody}{size=\normalsize}

\setbeamertemplate{footline}{
  \begin{beamercolorbox}[wd=\paperwidth]{upper separation line foot}
    \rule{0pt}{2pt}
  \end{beamercolorbox}

  \begin{beamercolorbox}{footline}%
    \vskip 0.7em
    \begin{columns}[t]
	    \begin{column}{0.025\linewidth}
	    \end{column}
	    \begin{column}{0.19\linewidth}
		    \center		    \includegraphics[height=3em]{Logo_FG_HD}
	    \end{column}

	    \begin{column}{0.19\linewidth}
		    \center		    \includegraphics[height=3em]{grid5000}
	    \end{column}

	    \begin{column}{0.19\linewidth}
		    \center	    \includegraphics[height=3em]{logo_loria_complet_couleur-crop}
	    \end{column}

	    \begin{column}{0.19\linewidth}
		    \center    \includegraphics[height=2.5em]{NancyU}
	    \end{column}
	    \begin{column}{0.19\linewidth}
		    \center		    \includegraphics[height=2.5em]{INRIA_SCIENTIFIQUE_FR_CMJN_CROP.pdf}
	    \end{column}

	    \begin{column}{0.025\linewidth}
	    \end{column}

    \end{columns}
    \vskip 1.2em
  \end{beamercolorbox}
}

  \begin{document}
  \begin{frame}{}
	  \vskip -0.5em
	  \begin{columns}[t]
	  \begin{column}{0.48\linewidth}

	\begin{block}{Objectifs}
		\begin{itemize}
			\item \textbf{Utiliser Grid'5000 comme infrastructure de test pour gLite}
					\bigskip
			\begin{itemize}
				\item Pouvoir reproduire les conditions expérimentales nécessaires à une expérience, mais difficiles à rencontrer en production (exemple~: comportement d'une application face à un crash de services)
					\bigskip
				\item Pouvoir remplacer des parties de l'infrastructure par d'autres versions de composants du middleware (test de nouvelles versions, test d'interopérabilité)

					\bigskip
				\item Pouvoir réaliser et reproduire des expériences dans un environnement stable (pas de variations entre des expériences $\leadsto$ pouvoir comparer les résultats)
					\bigskip
				\item Éviter de perturber ou de surcharger l'infrastructure de production avec des tâches de test
			\end{itemize}
	\end{itemize}
	\end{block} \vskip 0.5em

	\begin{block}{Grid'5000}
	\begin{columns}
	\begin{column}{0.65\columnwidth}
		\begin{itemize}
			\item \textbf{Plate-forme expérimentale pour la recherche sur les systèmes distribués} et le calcul parallèle à hautes performances
					\bigskip
			\item 1700 machines (7000 coeurs) dans 10 sites en France
					\bigskip
			\item Reconfiguration par les utilisateurs~: remplacement du système installé sur les noeuds avec Kadeploy, isolation réseau avec KaVLAN
		\end{itemize}
	\end{column}
	\begin{column}{0.33\columnwidth}
		\centerline{\includegraphics[width=0.9\columnwidth]{g5k}}
	\end{column}
\end{columns}
	\end{block}

	\vskip 0.5em

	\begin{block}{Infrastructure gLite déployée}
\begin{itemize}

	\item Une \textbf{VO} et son \textbf{VOMS (\textsl{Virtual Organization Membership Service})},
	annuaire des utilisateurs
					\bigskip

\item Plusieurs sites, composés de~:

		\bigskip
\begin{enumerate}

	\item Un \textbf{BDII (\textsl{Berkeley Database Information Index})}, annuaire
		des ressources disponibles sur chaque site

		\bigskip

	\item Un \textbf{CE (\textsl{Computing Element})}, service de soumission des
		tâches à un site de calcul donné

	\item Des \textbf{worker nodes} et un \textbf{batch scheduler}
		permettant d'y accéder. Le couple \textbf{Torque/Maui} a été
		utilisé

		\bigskip

	\item Une \textbf{UI (\textsl{User Interface})}, interface d'accès pour les
		utilisateurs

\end{enumerate}
\end{itemize}
	\end{block}
      \end{column}
      \begin{column}{0.48\linewidth}
\begin{block}{Outils développés}
	\begin{columns}[b]
	\begin{column}{0.75\linewidth}
	\begin{itemize}
		\item \textbf{Image Scientific Linux 5.5} minimale et générique (fonctionnant sur tous les clusters de Grid'5000) pour l'outil de déploiement Kadeploy
	\end{itemize}
	\end{column}
	\begin{column}{0.2\linewidth}
		\includegraphics[width=0.98\columnwidth]{sl-logo-800}
	\end{column}
	\end{columns}
		\bigskip
	\begin{itemize}
		\item \textbf{Scripts Ruby} documentés permettant une installation totalement automatisée de gLite à partir des dépôts RPM
		\bigskip
		\begin{itemize}
			\item Description de la plate-forme à déployer (VO, sites, clusters) dans un fichier de configuration
			\bigskip
			\item Création d'une autorité de certification pour générer et signer automatiquement les certificats des utilisateurs et des machines
		\bigskip
			\item Pré-remplissage du cache RPM sur les noeuds avec Kadeploy pour accélérer le déploiement (intégré à l'image)
		\end{itemize}
	\end{itemize}
	\vskip 0.5em
	\begin{center}
		\textbf{\url{https://github.com/sbadia/gdeploy/}}
	\end{center}
	\end{block} \vskip 0.8em


\begin{block}{Résultats}
	\textbf{Utilisation de Grid'5000 pour déployer le middleware gLite}
\begin{itemize}
	\item Déploiement sur \textbf{5 sites, 8 clusters et 441 machines}
		\bigskip
\item Installation de l'ensemble des machines sous Scientific Linux 5.5 minimale avec Kadeploy~: 10~minutes
		\bigskip
\item Configuration de l'infrastructure gLite~: \textbf{40~minutes}
\end{itemize}
	\end{block} \vskip 0.8em

	\begin{block}{Perspectives}
	\begin{itemize}
		\item \textbf{Amélioration du script de déploiement}
		\begin{itemize}
		\item Déploiement de plusieurs VO
		\item Déploiement d'autres services gLite~:\\ stockage, supervision
		\bigskip
			\end{itemize}
		\item \textbf{Collaborations}
		\bigskip
			\begin{itemize}
				\item \textbf{Expériences sur des évolutions de composants du middleware gLite}
		\bigskip
				\item \textbf{Expériences sur des outils interagissant avec le middleware gLite~: moteurs de workflows, gestionnaires de jobs pilotes, etc.}
					\begin{itemize}
						\item Simulation de pannes de services
						\item Injection de charge
						\item Soumission d'un grand nombre de tâches factices
					\end{itemize}
			\end{itemize}
	\end{itemize}
	\end{block}
\end{column}
\end{columns}

\vskip 0.5em

    \begin{block}{Processus de déploiement}
	    \center
\begin{tikzpicture}[node distance=1em]
\tikzstyle{every task} = [thick,fill=i6colorscheme3]
\tikzstyle{every sequence} = [thick]
\tikzstyle{every gateway} = [thick,fill=i6colorscheme3]
\tikzstyle{every event} = [thick,fill=i6colorscheme3]
\tikzstyle{end event} = [event,line width=5pt,fill=i6colorscheme3]

	\node[event] (start) {};
	\node[task,right=of start] (updglite) { Update SL and \\ add gLite RPM \\ repositories \\ on all nodes };
	\draw[sequence,->] (start) -- (updglite);
	\node[task,right=of updglite] (ca) { Create \\ CA };
	\draw[sequence,->] (updglite) -- (ca);
	\node[task,right=of ca] (vo) { Create VO and \\configure VOMS};
	\draw[sequence,->] (ca) -- (vo);
	\node[gateway,right=of vo] (pg) {+};
	\draw[sequence,->] (vo) -- (pg);
	\node[task,above right=of pg] (bdii) {Configure \\ BDII};
	\draw[sequence,->] (pg) |- (bdii);
	\node[gateway,right=of bdii] (pg2) {+};
	\draw[sequence,->] (bdii) -- (pg2);
	\node[task,right=of pg2] (wn) {Configure \\ worker nodes\\ and batch};
	\draw[sequence,->] (pg2) -- (wn);
	\node[align=center,below=of wn] (etc2) { {\Large \ldots } \\ \small (on each cluster)};
%	\node[below right=of pg2] (etc2-target) {};
	\draw[sequence,->,loosely dashed] (pg2) |- (etc2);
	\node[gateway,below right=of wn] (pg3) {+};
	\draw[sequence,->,loosely dashed] (etc2) -| (pg3);
	\draw[sequence,->] (wn) -| (pg3);
	\node[task,right=of pg3] (ce) {Configure \\ CE};
	\draw[sequence,->] (pg3) -- (ce);
	\node[task,right=of ce] (ui) {Configure \\ UI};
	\draw[sequence,->] (ce) -- (ui);
	\node[gateway,right=of ui] (pg4) {+};
	\draw[sequence,->] (ui) -- (pg4);
	\node[end event,right=of pg4] (end) {};
	\draw[sequence,->] (pg4) -- (end);


	\node[align=center,below=of etc2] (etc1) { {\Large \ldots } \small (on each site)};
	\draw[sequence,->,loosely dashed] (pg) |- (etc1);
	\draw[sequence,->,loosely dashed] (etc1) -| (pg4);
  \end{tikzpicture}
     \end{block}
  \end{frame}
\end{document}


%%%%%%%%%%%%%%%%%%%%%%%%%%%%%%%%%%%%%%%%%%%%%%%%%%%%%%%%%%%%%%%%%%%%%%%%%%%%%%%%%%%%%%%%%%%%%%%%%%%%
%%% Local Variables:
%%% mode: latex
%%% TeX-PDF-mode: t
